\chapter*{Introduction}

These notes originated as lecture notes for the Mathematical Analysis course of %% README
the Physics degree program at the University of Pisa during the academic years %% README
from 2017/18 to 2024/25. %% README

The notes (as well as the course they refer to) concern the analysis of %% README
functions of a real variable. The topics covered include numerical series and %% README
sequences, recursive sequences, differential calculus, and integral calculus. A %% README
minimal amount of functional analysis is conducted with the aim of considering, %% README
as a final topic, the study of ordinary differential equations. Complex numbers %% README
are introduced from the outset and are used where they can help provide a more %% README
unified and conceptually simpler view of the topics covered (particularly in the %% README
study of power series, in the definition of trigonometric functions, and in the %% README
resolution of linear differential equations). %% README

The notes are extensive, with no attempt at conciseness. The goal is to gather %% README
all those results that cannot always be presented in a detailed and rigorous %% README
manner during lectures. For example, we will find the construction of real %% README
numbers, equivalent definitions of the exponential function and trigonometric %% README
functions (with $e$ and $\pi$). We propose the proof of the fundamental theorem %% README
of algebra, Stirling's formula, Wallis's formula, and the irrationality of $e$ %% README
and $\pi$. A formal definition of Landau's symbols, small $o$ and big $O$, is %% README
proposed along with the relevant theorems for handling these expressions. The %% README
same is done for the symbol of indefinite integral. Algebraic results are %% README
discussed that justify algorithms for calculating the primitives of rational %% README
functions and for solving linear differential equations using the similarity %% README
method. %% README

The first chapter, dedicated to formal systems, the construction of numerical %% README
sets, and the definition of elementary functions, aims to provide a mathematical %% README
foundation for all those notions that the student should have already learned in %% README
the school curriculum. Some topics covered in the first chapter may be difficult %% README
to understand as they are very technical: they can nevertheless be skipped as %% README
they deal with concepts that should at least intuitively be already known. In %% README
particular, the axiomatic definition of the set of real numbers is probably very %% README
different from how it is usually presented in most textbooks introducing %% README
mathematical analysis: instead of assuming the field properties (and thus the a %% README
priori existence of the two operations of addition and multiplication), only the %% README
addition operation is considered while multiplication is constructed %% README
consequently. This approach is justified by the observation that the %% README
construction of multiplication is exactly the same construction that is usually %% README
used to construct the operation of exponentiation. Thus, the same abstract %% README
theorem (isomorphism of densely and continuously ordered groups) is used to %% README
construct multiplication, exponential functions, and trigonometric functions. %% README

Figures are not frequent, but many of them have a \emph{QR-code} (a square %% README
formed by a cloud of pixels) that allows access to the figure \emph{online} and %% README
to modify its parameters. 
On the page \url{https://paolini.github.io/AnalisiUno/}, 
all links to interactive figures
are provided.
\begin{comment}
Below on this page, you will find a list of links to interactive figures. %% README
\end{comment}
 %% README
\myqrcode{https://paolini.github.io/AnalisiUno/}{}{this web page}

These notes are made freely available both in PDF format and in source form %% README
\begin{comment}
    LaTeX. %% README
\end{comment}
\LaTeX{}. 
One way to support this project and keep it freely available for %% README
everyone is to purchase a printed copy. %% README

The material is constantly evolving and certainly contains errors and %% README
inconsistencies. Any suggestions or comments are welcome! %% README

\section*{Contributions}

I thank:
%
Giosuè Aiello,
Valerio Amico,
Francesco Antonacci,
Filippo Belfortini,
Rico Bellani,
Ruggero Bellettini,
Fabio Bensch,
Jacopo Bernardini,
Paolo Bernieri,
Riccardo Bin,
Elia Bonistalli,
Nicolò Bottiglioni,
Giovanni Bozzi,
Antonino Calderone,
Davide Campanella Galanti,
Alessandro Canzonieri,
Giulio Carlo,
Luca Casagrande,
Alessandro Casini,
Marco Catapano,
Augusto Cattafesta,
Giulio Cianti,
Luca Ciceri,
Tommaso Ceccotti,
Giorgio Ciocca,
Luca Ciucci,
Dario Cordaro,
Mattia Cozzani,
Francesco Debortoli,
Martino Dimartino,
Igor Di Tota,
Edoardo Fedi,
Leonardo Fontana,
Diego Fortini,
Davide Geria,
Irene Giambastiani,
Giorgio Gioffré,
Lorenzo Guglielmi,
Mattia Roberto Klepser,
Michele Pio Iallorenzi,
Irene Iorio,
Marco Labella,
Davide Labrosciano,
Andrea La Mendola,
Viviana Lippolis,
Gabriele Lo Gioco,
Jacopo Lombardi,
Mattia Luconi Trmbacchi,
Matteo Lupo,
Marco Malucchi,
Francesco Maria Manuello,
Giovanni Marcenaro,
Luigina Mazzone,
Roberto Menta,
Grazia Emilia Miccoli,
Luca Milanese,
Silvia Modenese,
Michele Monti,
Elena Morano,
Aurora Mugnai,
Francesco Nagni,
Alice Nicoletta,
Marco Nuti,
Samuele Paglia,
Chiara Parducci,
Ruben Pariente,
Silvio Davide Pascale,
Daniele Pavarini,
Guglielmo Pellitteri,
Paolo Pennoni,
Davide Perrone,
Giuditta Pevere,
Lorenzo Pierfederici,
Alessandro Rayan Ahmed,
Riccardo Maria Rini,
Mattia Ripepe,
Francesco Rodriguez,
Chiara Rustici,
Elisa Sabadini,
Vincenzo Santarsiero,
Gabriele Scarci,
Gabriele Sclafani,
Maria Antonella Secondo,
Carmen Selicato,
Francesco Sermi,
Gabriele Siffredi,
Irene Silvestro,
Alessio Squilloni,
Jacopo Taddei,
Antonio Tagliente,
Francesco Enrico Teofilo,
Federico Tettamanti,
Laura Toni,
Giacomo Trupiano,
Bianca Turini,
Ettore Ulivieri,
Francesco Vaselli,
Antoine Venturini,
Diego Vignocchi,
Matteo Vilucchio,
Piero Viscone,
Dariel Vllamasi
%
who have reported errors and corrections.

I thank my colleagues Vincenzo Tortorelli, Pietro Majer, Matteo Novaga, and
Carlo Carminati who have given me many valuable suggestions. Thanks also to
Davide Lombardo who suggested a simplification of the proof of
theorem~\ref{th:indipendenza_successioni_esponenziali}.

\newpage

\section*{Pronunciation of Foreign Letters}

\emph{English:}
\begin{center}
\begin{minipage}{3cm}
$J$ $j$ gei (long \emph{i}) \index{j}\\
$K$ $k$ cappa (key) \index{k}\\
$W$ $w$ double-u \index{w}\\
$X$ $x$ ex \index{x}\\
$Y$ $y$ ipsilon (Greek \emph{i}) \index{y}
\end{minipage}
\end{center}
%
\emph{Greek:}
\begin{center}
\begin{minipage}{3cm}
$A$ $\alpha$ alpha \index{$\alpha$} \\
$B$ $\beta$ beta \index{$\beta$}\\
$\Gamma$ $\gamma$ gamma \index{$\gamma$} \\
$\Delta$ $\delta$ delta \index{$\delta$} \\
$E$ $\eps$ epsilon \index{$\eps$} \\
$Z$ $\zeta$ zeta\footnotemark[1] \\
$H$ $\eta$ eta  \\
$\Theta$ $\theta$ theta \index{$\theta$}
\end{minipage}%
\begin{minipage}{3cm}
$I$ $\iota$ iota  \\
$K$ $\kappa$ kappa\footnotemark[2]  \\
$\Lambda$ $\lambda$ lambda \index{$\lambda$} \\
$M$ $\mu$ mu \index{$\mu$} \\
$N$ $\nu$ nu \index{$\nu$} \\
$\Xi$ $\xi$ xi  \\
$O$ $o$ omicron  \\
$\Pi$ $\pi$ pi\footnotemark[3] \index{$\pi$}
\end{minipage}%
%
\begin{minipage}{3cm}
$P$ $\rho$ rho \index{$\rho$} \\
$\Sigma$ $\sigma$ sigma \index{$\sigma$} \index{$\Sigma$} \\
$T$ $\tau$ tau  \index{$\tau$}\\
$Y$ $\upsilon$ upsilon\footnotemark[4]\\
$\Phi$ $\phi$ phi \index{$\phi$} \\
$X$ $\chi$ chi \index{$\chi$} \\
$\Psi$ $\psi$ psi \index{$\psi$} \\
$\Omega$ $\omega$ omega \index{$\omega$}
\end{minipage}
\end{center}
%
\emph{Hebrew:}
\begin{center}
\begin{minipage}{3cm}
$\aleph$ aleph
\end{minipage}
\end{center}

\footnotetext[1]{usually pronounced with a soft $z$ to distinguish it from $z$.}
\footnotetext[2]{in English, the pronunciation between $k$ \emph{key} and $\kappa$ \emph{kappa} is distinguished.}
\footnotetext[3]{in Italian \emph{pi greco}, in English \emph{pai}, to distinguish from $p$.}
\footnotetext[4]{pronounced \emph{upsilon} to distinguish it from $y$.}